\chapter{Trader}\label{ch:in:trader}

Le 4 mai 2021, le plus grand groupe boursier américain de dérivés annonçait qu'il ne rouvrirait pas les corbeilles de négociation qu'il avait
fermées pendant la pandémie, à l'exception d'une seule corbeille d'options~; en 2025, 93~\% de ses volumes étaient électroniques. Le métier de trader avait
changé des cris aux clics, puis des clics à la supervision d'algorithmes qui cotent et négocient plus vite que n'importe quelle personne.
Ce qui reste au trader, c'est le jugement~: quels risques détenir, comment régler les machines, quand les arrêter et comment
en répondre. Ce chapitre ouvre la partie~III, un chapitre par métier, et chacun suit le même plan~: ce qu'est le métier chez chaque type d'employeur
comment il a évolué, quels contrôles et quelles règles l'encadrent, une journée documentée ou composite, les données de rémunération
du chapitre~14 pour ce métier, et une analyse qui montre en quoi consiste le travail du métier. Pour le trader, l'analyse porte sur une
question que tout desk affronte~: combien d'algorithmes une personne peut-elle surveiller~?

\begin{center}\footnotesize
\begin{tabular}{@{}p{3.2cm}p{9.4cm}@{}}
\toprule
\multicolumn{2}{@{}l}{\textbf{Fiche métier~: trader}}\\
\midrule
employeurs & teneurs de marché, banques, hedge funds, plateformes\\
horizon des décisions & de quelques secondes à quelques jours\\
P\&L & détient un portefeuille (\omterm{def:in:trader:kinds}{trader de risque}) ou aucun (\omterm{def:in:trader:kinds}{trader d'exécution})\\
rend compte à & responsable de desk\\
codes professionnels dans les dépôts & 13-2099.01, 13-2051, 41-3031~; intitulés contenant «~trader~» ou «~trading analyst~»\\
où la série l'enseigne & livre~1, chapitre~2~; livre~10, chapitre~24~; livre~11, chapitre~27\\
\bottomrule
\end{tabular}
\end{center}

\section{Trois traders~: teneur de marché, desk de flux bancaire, hedge fund}

Le mot «~trader~» recouvre des métiers qui ne partagent guère plus qu'un écran. Deux définitions séparent les principales catégories.

\begin{definition}[Trader de risque, trader d'exécution]\label{def:in:trader:kinds}
Un \emph{trader de risque}\index{trader de risque} (risk trader) décide quelles positions la société détient, dans les limites fixées, et porte le
résultat d'un portefeuille. Un \emph{trader d'exécution}\index{trader d'exécution} (execution trader) exécute des ordres décidés par quelqu'un d'autre, un gérant
de portefeuille ou un client, et son évaluation porte sur le coût et la qualité d'exécution plutôt que sur le profit de la position.
\end{definition}

\begin{itemize}
\item \textbf{Chez un teneur de marché}, le trader est un \omterm{def:in:trader:kinds}{trader de risque} dont le portefeuille est l'inventaire que lui laissent
ses cotations (livre~11, chapitres~1 à 3). Un teneur de marché décrit ses traders quantitatifs comme «~au cœur de l'action~: ils suivent
les marchés de près, prenant des décisions rapides~», et ajoute qu'«~ils développent des outils pour automatiser les décisions de trading~»~; ses
traders institutionnels «~interagissent avec les courtiers et les autres équipes de trading~».
\item \textbf{Au desk de flux d'une banque}, le trader fixe des prix aux clients de la banque par l'intermédiaire de la vente (livre~1, chapitre~2) et
gère le risque qui en résulte~: un \omterm{def:in:trader:kinds}{trader de risque} dont le portefeuille est le flux de la franchise (livre~9, chapitre~24).
\item \textbf{Chez un hedge fund}, le trader est en général un \omterm{def:in:trader:kinds}{trader d'exécution} au service des gérants de portefeuille, qui
passe les ordres par l'intermédiaire de courtiers et d'algorithmes (livre~10, chapitre~24)~; certains fonds confient aussi à leurs traders un risque propre.
\end{itemize}
Les trois partagent la discipline des limites et l'obligation de savoir, à chaque instant, ce que la société détient~; ils diffèrent
par ceux qu'ils servent, ce qu'ils détiennent et la manière dont leur rémunération est évaluée (chapitre~13).

\section{Du click trading à la supervision d'algorithmes}

\begin{definition}[Click trader, superviseur d'algorithmes]\label{def:in:trader:supervisor}
Un \emph{click trader}\index{click trader} saisit et gère des ordres à la main sur un écran électronique.
Un \emph{superviseur d'algorithmes}\index{superviseur d'algorithmes} (algorithm supervisor) est un trader chargé d'un ou de plusieurs algorithmes de trading, qui
en règle les paramètres, les surveille en temps réel, traite leurs alertes et les arrête lorsqu'ils se comportent mal.
\end{definition}

Le trader à la corbeille du chapitre~2 lisait la foule~; le \omterm{def:in:trader:supervisor}{click trader} lisait l'écran~; le \omterm{def:in:trader:supervisor}{superviseur d'algorithmes} lit les
algorithmes. Le changement n'était pas seulement technique. Les règles européennes sur le trading algorithmique font du superviseur une fonction nommément désignée~: la
surveillance en temps réel du trading algorithmique «~est assurée par le trader chargé de l'algorithme de trading ou
de la stratégie de trading algorithmique, ainsi que par la fonction de gestion des risques~», et le personnel responsable «~répond
aux problèmes opérationnels et réglementaires dans des délais appropriés~». La journée d'un superviseur est donc une file d'attente~: des alertes arrivent
des algorithmes~: chacune demande une intervention, et une alerte qui attend est un risque qui court sans surveillance. Combien d'algorithmes
une personne peut superviser est une question de files d'attente, et le chapitre y répond dans la section~5.

\section{Contrôles et responsabilité~: enregistrements, certification, kill switches}

\begin{definition}[Régime de certification]\label{def:in:trader:certification}
Un \emph{régime de certification}\index{régime de certification} est une règle qui oblige une société à évaluer, à l'embauche et au moins
chaque année, que les personnes occupant certaines fonctions nommées sont aptes et honorables pour ces fonctions, et à certifier cette évaluation~; au Royaume-Uni, les fonctions concernées comprennent
les traders pour compte propre susceptibles de causer un préjudice important, ainsi que les personnes qui approuvent ou surveillent les algorithmes de trading.
\end{definition}

\begin{dated}{2026-07}{Les règles qui s'imposent au trader}\label{dat:in:trader:controls}
\textbf{États-Unis}~: un trader de titres chez un broker-dealer qui négocie pour son propre compte hors d'une bourse s'enregistre comme
Securities Trader Representative, après les examens Securities Industry Essentials et Series~57. \textbf{Royaume-Uni}~: agir comme
«~un trader pour compte propre dont l'activité implique, ou pourrait impliquer, un risque de préjudice important pour la société ou pour
l'un de ses clients~» constitue une fonction de certification, de même qu'approuver le déploiement d'un algorithme de trading et avoir
une responsabilité importante dans la surveillance de la conformité d'un algorithme (FCA Handbook, SYSC 27.8). \textbf{Union européenne}~: une société qui négocie
algorithmiquement «~doit pouvoir annuler immédiatement, à titre de mesure d'urgence, tout ou partie de ses ordres non exécutés
~» (la fonctionnalité de kill), et doit surveiller ses algorithmes en temps réel par l'intermédiaire du trader responsable et de sa
fonction de gestion des risques (Delegated Regulation 2017/589, articles 12 et 16).
\end{dated}

Les règles transforment le jugement du trader en une responsabilité documentée. Un kill switch (livre~11, chapitre~27) existe au niveau de la société et
de la plateforme~; une limite de perte arrête un portefeuille lorsque ses pertes atteignent un montant fixé~; un contrôle de risque pré-négociation rejette un
ordre qui enfreint une limite avant qu'il ne quitte la société. La part du trader consiste à savoir lesquels de ces contrôles s'appliquent, à
régler les limites que la société lui délègue, et à actionner l'arrêt lorsqu'il est nécessaire, et non lorsqu'il est trop tard. Le chapitre~24 décrit
les fonctions de risque et de conformité, de l'autre côté de ces contrôles.

\section{Une journée documentée}

Les règles et la description publiée par une société permettent de reconstituer une journée composite d'un \omterm{def:in:trader:supervisor}{superviseur d'algorithmes} chez un teneur de marché~;
c'est un composite, et non l'emploi du temps d'une société en particulier.
\begin{enumerate}
\item \textbf{Avant l'ouverture}~: vérifier les paramètres, les limites et les chemins de kill des algorithmes~; passer en revue les événements de la nuit et
les positions~; convenir du plan de la journée avec les chercheurs et les développeurs.
\item \textbf{L'ouverture}~: les minutes les plus chargées, lorsque les cotations s'élargissent et que les alertes arrivent le plus vite~; le superviseur suit les exécutions, les
positions et les signaux de santé propres aux algorithmes.
\item \textbf{La séance}~: ajuster les paramètres à mesure que les conditions changent, répondre aux alertes, enquêter sur tout élément inexpliqué et
transmettre au risque et à la technologie tout problème qui dépasse le desk.
\item \textbf{Incidents}~: arrêter un algorithme, un produit ou l'ensemble du portefeuille lorsque le comportement n'est pas compris, et consigner les raisons de cet arrêt.
\item \textbf{Après la clôture}~: rapprocher les positions et le résultat, passer en revue les alertes et les exécutions de la journée, et réinjecter ce qui
a été appris dans les modèles et les contrôles.
\end{enumerate}
La journée du trader de flux bancaire ajoute les clients~: les prix demandés par la vente, les transactions à couvrir et une franchise à protéger
(chapitre~23). La journée du \omterm{def:in:trader:kinds}{trader d'exécution} d'un hedge fund suit les ordres des gérants de portefeuille et les algorithmes des courtiers.

\begin{dated}{2025-09}{Ce que gagnent les traders~: les données publiques}\label{dat:in:trader:pay}
Demandes de certification des conditions de travail aux États-Unis, exercice fiscal 2025, famille de métiers «~trader~»~: \omterm{def:in:how-pay-works:base}{salaire de base} médian proposé de 195\,000~\$ chez les teneurs de marché
(65 demandes, 9 employeurs~; écart interquartile de 150\,000~\$ à 250\,000~\$) et de 235\,000~\$ dans les banques (58 demandes, 5
employeurs~; de 152\,500~\$ à 282\,500~\$)~; les cellules des fonds systématiques et des plateformes sont supprimées. L'enquête professionnelle publique du gouvernement
pour les agents commerciaux en valeurs mobilières, matières premières et services financiers du secteur des valeurs mobilières, mai 2025~: médiane
103\,030~\$, du 10\up{e}~centile au 90\up{e}~centile, de 55\,130~\$ à 309\,440~\$. \omterm{def:in:how-pay-works:base}{Salaire de base} seulement dans la première, la plupart des primes étant exclues de la
seconde (chapitre~14).
\end{dated}

\begin{omfigure}
\begin{tikzpicture}
\begin{axis}[width=10.5cm,height=4.6cm,xmin=0,xmax=340,ymin=0.4,ymax=3.6,ytick=data,
  yticklabels={{OEWS, agents commerciaux en valeurs mobilières},{LCA, traders en banque},{LCA, traders chez les teneurs de marché}},yticklabel style={font=\scriptsize},
  xticklabel style={font=\footnotesize},xlabel={\small milliers de \$ par an},grid=major,grid style={gray!20},
  table/col sep=comma]
\addplot[draw=none,mark=none,error bars/.cd,x dir=both,x explicit,error mark=none,error bar style={omDef,line width=1pt}]
  table[x=p50,y=pos,x error plus expr=\thisrow{p90}-\thisrow{p50},x error minus expr=\thisrow{p50}-\thisrow{p10}]
  {figdata/industry/16-trader/pay.csv};
\addplot[draw=none,mark=none,error bars/.cd,x dir=both,x explicit,error mark=none,error bar style={omDef!40,line width=7pt}]
  table[x=p50,y=pos,x error plus expr=\thisrow{p75}-\thisrow{p50},x error minus expr=\thisrow{p50}-\thisrow{p25}]
  {figdata/industry/16-trader/pay.csv};
\addplot[only marks,mark=|,mark size=5pt,line width=1.5pt,omThm] table[x=p50,y=pos]{figdata/industry/16-trader/pay.csv};
\end{axis}
\end{tikzpicture}

\omcaption{Éléments de rémunération des traders~: le 10\up{e}~centile au 90\up{e}~centile (fin), l'écart interquartile (épais) et la médiane
(repère) du \omterm{def:in:how-pay-works:base}{salaire de base} proposé dans les demandes de certification de l'exercice fiscal 2025, à côté des salaires de l'enquête professionnelle des agents commerciaux en valeurs mobilières du secteur des valeurs mobilières,
mai 2025. Données~: \texttt{data/industry/lca\_ranges.csv},
\texttt{oews\_finance.csv}, via \texttt{in\_trader.pay}.}\label{fig:in:trader:pay}
\end{omfigure}

Les dépôts font apparaître des \omterm{def:in:how-pay-works:base}{salaires de base} compris entre 150\,000~\$ et 250\,000~\$ pour la plupart des traders parrainés, dans les deux types d'employeurs,
avec une médiane plus élevée dans les banques~; l'enquête regroupe sous la même profession des traders, des courtiers et des conseillers, ce qui explique que sa médiane
soit la moitié de celle des dépôts. La prime, que le chapitre~14 a trouvée capable de multiplier le \omterm{def:in:how-pay-works:base}{salaire de base} plusieurs fois chez un teneur de marché, n'est incluse dans aucune des deux.

\section{Tutoriel~: un superviseur, de nombreux algorithmes}\label{tut:in:trader:tut}

\begin{tutorial}
\textbf{Objectif.} Trouver combien d'algorithmes un superviseur peut surveiller sans que les alertes attendent, avec des taux donnés.
\textbf{État final~:} \cref{fig:in:trader:pay,fig:in:trader:tail,fig:in:trader:capacity}.
\begin{tutsteps}
\item \textbf{Éléments de rémunération.} \texttt{firm.roles.lca\_cells(rows, role)} avec le rôle \texttt{trader} lit les plages du chapitre~14 pour la famille de métiers,
avec ses cellules supprimées~; la ligne de l'enquête professionnelle provient de \texttt{oews\_finance.csv}.
\item \textbf{La file d'attente.} Chaque algorithme déclenche des alertes à un rythme illustratif de 1 par heure, selon un flux de Poisson~; le traitement dure
30 secondes en moyenne~; la norme veut qu'au plus 1~\% des alertes attendent plus d'une minute. Avec un traitement exponentiel,
le temps d'attente dépasse $t$ avec la probabilité $\rho e^{-(\mu-\lambda)t}$ (livre~4, chapitre~8).
\omcode{code/firm/roles/firm_roles.py}{79}{93}{La queue du temps d'attente M/M/1 et le plus grand nombre d'algorithmes
qu'un superviseur peut surveiller dans le cadre de la norme.}\label{lst:in:trader:mm1}
\item \textbf{Une queue plus lourde.} \texttt{lindley\_tail} simule la file avec des durées de traitement lognormales (moyenne 30 secondes,
écart-type 45 secondes) par la récurrence de Lindley.
\end{tutsteps}
\end{tutorial}

Avec ces taux, le superviseur respecte la norme pour 7 algorithmes~: à 7 algorithmes, 0.89~\% des alertes attendent plus
d'une minute, et à 8, 1.03~\%. La réponse est sensible au taux d'alertes~: à une demi-alerte par heure, elle vaut 15~; à deux alertes, 3~; à
quatre, seulement 1 (\cref{fig:in:trader:capacity}). Elle l'est encore davantage à la norme~: admettre cinq minutes au lieu
d'une minute fait monter le nombre à 71 pour une alerte par heure. Et la queue de la durée de traitement compte~: avec un traitement lognormal
de même moyenne, 1.37~\% des alertes attendent plus d'une minute à 7 algorithmes (\cref{fig:in:trader:tail}).

\begin{omfigure}
\begin{tikzpicture}
\begin{axis}[width=10.5cm,height=5.4cm,xmin=0.5,xmax=20.5,ymin=0,ymax=4.5,xlabel={\small algorithmes supervisés},
  ylabel={\small alertes en attente depuis plus d'une minute, \%},tick label style={font=\footnotesize},grid=major,grid style={gray!20},
  table/col sep=comma,legend style={legend cell align=left,font=\scriptsize,at={(0.03,0.97)},anchor=north west}]
\addplot[omDef,thick,mark=*,mark size=1.5pt] table[x=n,y=mm1]{figdata/industry/16-trader/tails.csv};
\addplot[omThm,thick,mark=square*,mark size=1.5pt] table[x=n,y=lognormal]{figdata/industry/16-trader/tails.csv};
\addplot[black!60,dashed] coordinates {(0.5,1) (20.5,1)};
\legend{traitement exponentiel (forme fermée),traitement lognormal (simulé),la norme de 1~\%}
\end{axis}
\end{tikzpicture}

\omcaption{Part des alertes qui attendent plus d'une minute en fonction du nombre d'algorithmes, chacun déclenchant une alerte par
heure, avec un traitement de moyenne 30 secondes. Les taux sont illustratifs. Données~: \texttt{in\_trader.tails} (400\,000 alertes simulées
par point).}\label{fig:in:trader:tail}
\end{omfigure}

\begin{omfigure}
\begin{tikzpicture}
\begin{axis}[width=10.5cm,height=5.2cm,ybar,area legend,bar width=12pt,ymode=log,log origin=infty,ymin=0.8,ymax=300,xmin=0.4,xmax=4.6,
  xtick={1,2,3,4},xticklabels={0.5,1,2,4},log ticks with fixed point,
  xlabel={\small alertes par algorithme et par heure},ylabel={\small algorithmes par superviseur},
  tick label style={font=\footnotesize},grid=major,grid style={gray!20},table/col sep=comma,
  legend style={legend cell align=left,font=\scriptsize,at={(0.5,-0.3)},anchor=north,/tikz/every even column/.append style={column sep=8pt}},
  legend columns=2,point meta=rawy,nodes near coords,nodes near coords style={font=\scriptsize}]
\addplot[fill=omDef!55,draw=omDef] table[x=pos,y=one_minute]{figdata/industry/16-trader/capacity.csv};
\addplot[fill=omThm!45,draw=omThm] table[x=pos,y=five_minutes]{figdata/industry/16-trader/capacity.csv};
\legend{1~\% attendent plus d'une minute,1~\% attendent plus de cinq minutes}
\end{axis}
\end{tikzpicture}

\omcaption{Nombre d'algorithmes qu'un superviseur peut surveiller, selon le taux d'alertes, pour deux normes de réponse (traitement exponentiel
de moyenne 30 secondes~; échelle logarithmique). Les taux sont illustratifs. Données~: \texttt{in\_trader.capacity}.}\label{fig:in:trader:capacity}
\end{omfigure}

\begin{example}[Mise en commun]\label{ex:in:trader:pool}
Un desk fait tourner 40 algorithmes à raison d'une alerte par heure chacun. Répartis en carnets distincts de 7 au plus, ils exigent 6 superviseurs pour
respecter la norme d'une minute. Avec la même norme et une file commune, 2 superviseurs suffisent~: à 40 alertes par heure,
seulement 0.17~\% des alertes attendent plus d'une minute, car une alerte n'attend que si les deux superviseurs sont occupés~; deux superviseurs mis en commun
peuvent prendre en charge 76 algorithmes, et trois, 174. Le prix de la mise en commun est que chaque superviseur doit pouvoir traiter
l'alerte de n'importe quel algorithme.
\end{example}

Pour un desk, la leçon n'est pas un chiffre mais une conception~: moins d'alertes, mais de meilleure qualité, augmentent la portée d'un superviseur plus que toute autre
modification, ce qui explique qu'une grande part du travail d'un trader superviseur, et de celui des développeurs à ses côtés (chapitre~19), porte
sur la décision de ce qui mérite une alerte.

\textbf{Pour aller plus loin.} Donner deux priorités aux alertes et traiter d'abord les urgentes~; laisser les taux d'alertes monter avec la volatilité des marchés
(chapitre~12) et trouver le nombre de superviseurs dont un desk a besoin lors de ses journées les plus chargées.

\section{Construction~: le registre des rôles et le modèle de supervision}\label{bld:in:trader:roles}

\begin{build}
\bfield{Rôle} Enregistrer chaque rôle de la partie~III sous forme de données, et tenir une analyse par rôle~; ce premier incrément ajoute le
registre, la fiche du trader et le modèle de supervision.
\bfield{Interface} \texttt{firm.roles}~:
\texttt{RoleCard(name, firm\_types, horizon, pnl, reports\_to, soc\_codes, title\_patterns, books, chapter)}~; \texttt{REGISTRY}, \texttt{register}, \texttt{card}~;
\texttt{lca\_cells(rows, role, fy, level)}~; \texttt{mm1\_wait\_tail(lam, mu, t)}~; \texttt{max\_strategies(a, h, t, p)}~;
\texttt{lindley\_tail(lam, service, t, n, rng)}. NumPy.
\bfield{Règles} Un rôle est enregistré une seule fois~; ses cellules de rémunération proviennent du tableau du chapitre~14, la suppression étant conservée~; les
taux du modèle de supervision sont ceux de l'appelant, étiquetés illustratifs.
\bfield{Tests d'acceptation} \texttt{code/firm/roles/tests/}~: l'enregistrement d'un rôle deux fois échoue~; la queue M/M/1 vaut 1 à
pleine charge ou au-delà et concorde avec la file simulée à l'erreur d'échantillonnage près~; le plus grand nombre d'algorithmes est croissant
avec le taux d'alertes et avec la norme.
\bfield{Pour aller plus loin} Priorités et plusieurs superviseurs (M/M/c)~; alertes qui se regroupent en rafales~; traitement qui dépend de la
charge du superviseur.
\end{build}

\begin{omsources}
\item Communiqué de presse du CME Group du 4 mai 2021~; Form 10-K pour 2025.
\item FINRA, Series 57~; FCA Handbook, SYSC 27.8~; Commission Delegated Regulation (EU) 2017/589, articles 12 et 16.
\item Optiver, «~Breaking down the trading industry~».
\item Tableaux du chapitre~14 issus des fichiers LCA du Department of Labor et de l'enquête professionnelle du BLS.
\end{omsources}

\section{Exercices}

\begin{exercise}[$\star$]\label{exo:in:trader:1}
Classer en \omterm{def:in:trader:kinds}{trader de risque} ou \omterm{def:in:trader:kinds}{trader d'exécution}~: le trader d'options d'un teneur de marché~; un trader de hedge fund qui exécute l'ordre d'un gérant de portefeuille~;
le trader de taux d'une banque qui price un swap client et le couvre.
\end{exercise}

\begin{exercise}[$\star$]\label{exo:in:trader:2}
Dans une file M/M/1 avec $\lambda=7$ alertes par heure et $\mu=120$, calculer la probabilité qu'une alerte attende plus d'une
minute.
\end{exercise}

\begin{exercise}[$\star$]\label{exo:in:trader:3}
Quelles trois règles de la boîte datée s'appliqueraient à un \omterm{def:in:trader:supervisor}{superviseur d'algorithmes} dans une société britannique négociant sur des plateformes de l'UE, et
ce que chacune exige de la personne~?
\end{exercise}

\begin{exercise}[$\star\star$]\label{exo:in:trader:4}
Montrer qu'à norme fixée, le nombre d'algorithmes varie à peu près en proportion inverse du taux d'alertes lorsque le
superviseur est peu chargé.
\end{exercise}

\begin{exercise}[$\star\star$]\label{exo:in:trader:5}
Pourquoi la médiane de l'enquête professionnelle pour les agents commerciaux en valeurs mobilières vaut-elle environ la moitié de la médiane des déclarations de main-d'œuvre pour les traders~?
\end{exercise}

\begin{exercise}[$\star\star$]\label{exo:in:trader:6}
À temps de traitement moyen égal, pourquoi un temps de traitement lognormal d'écart-type 45 secondes rend-il les attentes
plus longues qu'avec un temps exponentiel~?
\end{exercise}

\begin{exercise}[$\star\star\star$]\label{exo:in:trader:7}
\emph{Programmation.} Avec deux superviseurs partageant une même file (M/M/2), combien d'algorithmes respectent la norme d'une minute à une
alerte par heure~? Simuler ou utiliser la formule d'Erlang~C.
\end{exercise}

\begin{exercise}[$\star\star\star$]\label{exo:in:trader:8}
\emph{Trouver la faille.} «~Nos superviseurs traitent chaque alerte en 30 secondes et il y a 60 alertes par heure et par personne, donc
ils ne sont occupés qu'à moitié et pourraient surveiller deux fois plus d'algorithmes.~»
\end{exercise}

\section{Problème~: un trader, de nombreux algorithmes}

\begin{problem}[{Problème du week-end --- un trader, de nombreux algorithmes}]\label{pb:in:trader:1}
Un teneur de marché fait passer son desk d'options du click trading à des algorithmes supervisés et demande combien d'algorithmes chaque
trader devrait surveiller.

\medskip\noindent\textbf{Partie I --- Le rôle.}
\begin{enumerate}
\item Définir \omterm{def:in:trader:kinds}{trader de risque} et \omterm{def:in:trader:kinds}{trader d'exécution}~; donner un exemple de chacun.
\item Décrire le trader chez un teneur de marché, un desk de flux d'une banque et un hedge fund.
\item Définir \omterm{def:in:trader:supervisor}{click trader} et \omterm{def:in:trader:supervisor}{superviseur d'algorithmes}.
\item Quelle décision le plus grand groupe américain de bourses de dérivés a-t-il prise en 2021, et quelle part de son volume était électronique en
2025~?
\item Qu'exigent les règles de l'UE du trader chargé d'un algorithme~?
\end{enumerate}

\medskip\noindent\textbf{Partie II --- Contrôles.}
\begin{enumerate}[resume]
\item Définir un \omterm{def:in:trader:certification}{régime de certification} et nommer les fonctions britanniques qui concernent les traders.
\item Quelle immatriculation un trader de titres américain chez un broker-dealer doit-il obtenir~?
\item Qu'est-ce que la fonctionnalité de kill switch, et qui doit y avoir accès~?
\item Décrire une journée type d'un \omterm{def:in:trader:supervisor}{superviseur d'algorithmes}.
\item Que disent les éléments de rémunération, et que laissent-ils de côté~?
\end{enumerate}

\medskip\noindent\textbf{Partie III --- La file.}
\begin{enumerate}[resume]
\item Énoncer le modèle~: arrivées, traitement, norme.
\item Écrire la queue du temps d'attente de la M/M/1 et la calculer pour 7 puis 8 algorithmes.
\item Comment la réponse change-t-elle avec le taux d'alertes~?
\item Comment change-t-elle avec une norme de cinq minutes~?
\item Comment change-t-elle avec un traitement lognormal~?
\end{enumerate}

\medskip\noindent\textbf{Partie IV --- Le verdict.}
\begin{enumerate}[resume]
\item Énoncer le \emph{résultat nommé}~: le plus grand nombre d'algorithmes qu'un superviseur peut encadrer avec au plus 1~\% de chances
qu'une alerte attende plus d'une minute, et sa sensibilité au taux d'alertes.
\item Quelle modification accroît le plus la portée d'un superviseur~?
\item Pourquoi le taux d'utilisation est-il un mauvais guide de la capacité ici~?
\item Que doit mesurer le desk avant de choisir le nombre~?
\item En deux phrases, répondre au teneur de marché.
\end{enumerate}
\end{problem}

\section{Questions d'entretien}

\begin{interviewq}[$\star$ \iqroles{trader}]\label{iq:in:trader:1}
Votre algorithme commence à perdre de l'argent d'une manière que vous ne comprenez pas. Que faites-vous dans les soixante secondes qui suivent~?
\end{interviewq}

\begin{interviewq}[$\star$ \iqroles{trader, risque}]\label{iq:in:trader:2}
Quelle est la différence entre un kill switch et une limite de perte~?
\end{interviewq}

\begin{interviewq}[$\star\star$ \iqroles{trader, chercheur}]\label{iq:in:trader:3}
Les alertes arrivent à raison de 60 par heure et chacune prend 30 secondes. Quelle fraction du temps le superviseur est-il occupé, et quel est
le temps d'attente moyen dans une file M/M/1~?
\end{interviewq}

\begin{interviewq}[$\star\star$ \iqroles{développeur, trader}]\label{iq:in:trader:4}
Concevez les alertes d'un algorithme de market making sur options. Lesquelles déclencheriez-vous en astreinte, et lesquelles vous limiteriez-vous à journaliser~?
\end{interviewq}

\begin{interviewq}[$\star\star$ \iqroles{banque, trader}]\label{iq:in:trader:5}
Un client demande un prix pour une taille importante juste avant une publication de données. Comment décidez-vous du prix et de la couverture~?
\end{interviewq}

\begin{interviewq}[$\star\star\star$ \iqroles{trader, chercheur}]\label{iq:in:trader:6}
Comment décideriez-vous si un trader apporte une valeur ajoutée à un ensemble d'algorithmes, plutôt que de se contenter de les surveiller~?
\end{interviewq}
